\section{Discussion}
\label{sec:discussion}

\subsection{Why v1 was silent and v2 is not}

The v1$\to$v2 contrast isolates two factors that together turn a silent mechanism into a measurable one. They are independently small but jointly load-bearing.

\paragraph{Floor.} Without a floor, the optimizer has a degree of freedom it does not need: the gate parameter $\alpha$ can drift toward $-\infty$ at no cost to the loss, because the rest of the network adapts to the resulting (zero) feedback. From an optimization standpoint, this is the path of least resistance --- it removes a noisy signal during training. The floor reparameterization (Equation~\ref{eq:gate-floor}) eliminates this degree of freedom: the optimizer can still learn an attenuation pattern, but cannot drive the contribution to zero. The mechanism is structurally guaranteed to remain in the loop. Empirically, even with the floor, the gate \emph{does} attenuate over training (from $0.55 \to 0.39$, Figure~\ref{fig:gate}), suggesting the optimizer would prefer it lower; the floor's role is to prevent it from going all the way down.

\paragraph{Mask.} The level mask focuses the cross-attention on the levels where small-object information lives. A single set of attention weights cannot simultaneously be optimal for stride-4 tokens (where small objects are sub-pixel-precise) and stride-32 tokens (where each token covers a $32\!\times\!32$ neighborhood). v1 forced the same mechanism to do both; v2 specializes. This is the same locality-prior that motivates feature pyramids in the first place~\cite{lin2017fpn} --- we extend it to the feedback path.

These two changes are complementary: the floor keeps the gate alive, and the mask makes the alive-gate's signal targeted. Stripping either change reverts the mechanism toward silence. We did not run the targeted ablation that would isolate floor-only vs. mask-only because we ran out of compute budget; this is one of the natural follow-ups (Section~\ref{sec:limitations}).

\subsection{What the $+0.99$ delta means and does not mean}

The $\Delta\mathrm{AP}_S = +0.99$ is the difference between two evaluations of the same trained network, differing only in whether one cross-attention block fires at inference. It is therefore an estimate of the \emph{causal contribution of the feedback computation}, conditional on the network's other weights having been trained \emph{with} the feedback active. The number is not a comparison of two trained networks. It does not say "training with feedback gives $+0.99$ AP$_S$ over training without" --- that would require a separately-trained no-feedback baseline that we did not run. It is the inference-time reading of how much the feedback module is doing, given that everything else has co-adapted to it during training.

This interpretation aligns with the negative result for v1: v1's mechanism was \emph{trained} (the cross-attention weights are non-trivial), but at inference the gate had decayed to ${\approx}0.12$ and the cross-attention output was multiplied by that small number before being added to memory. The result was statistically identical to skipping the cross-attention entirely. v1's failure is therefore a runtime failure, not a training failure.

\subsection{Limitations}
\label{sec:limitations}

We acknowledge several limitations of this study, in roughly decreasing order of how much they would change the conclusion:

\begin{enumerate}
    \item \textbf{Single seed.} All numbers in Tables~\ref{tab:main} and~\ref{tab:approaches} are from one training run. A multi-seed sweep would tighten the confidence interval on $\Delta\mathrm{AP}_S$ and check whether $+0.99$ is reproducible.
    \item \textbf{No floor-vs-mask isolation.} v2 changes both the gate parameterization and the level mask simultaneously, so we cannot say from this experiment alone whether one of the two is sufficient or whether they are jointly necessary. A 2$\times$2 ablation would resolve this.
    \item \textbf{Absolute $\mathrm{AP}_S$ is below the RT-DETR paper baseline.} v2 ON @ 640 reaches $\mathrm{AP}_S = 34.25$ vs. the published $34.7$. This is because we finetune from the public weights for 13 epochs rather than train from scratch on the original schedule, and we lack the multi-seed variance estimate to claim parity. Our claim is the $+0.99$ causal contribution, not a new state of the art.
    \item \textbf{Latency cost is dominant from P2, not feedback.} The $\sim$2$\times$ slowdown over the RT-DETR baseline is essentially entirely from adding the stride-4 feature level. The feedback module itself, especially with the level mask, is a small fraction of inference cost. Future work could explore whether the feedback approach generalizes to the original 3-level (P3--P5) configuration at full speed.
    \item \textbf{No cross-dataset evaluation.} Our results are on COCO only. Small-object detection has natural test domains (aerial imagery, traffic surveillance) where the gain might be larger or smaller; we did not run those.
\end{enumerate}

\subsection{Connections to prior work}

Iterative refinement of detection predictions is a long-standing idea in computer vision (e.g., cascade R-CNN~\cite{cai2018cascade}, deformable refinement layers in DETR variants~\cite{zhu2021deformable,liu2022dabdetr}). What we add is iterative refinement of the \emph{encoder memory} based on early decoder predictions, rather than iterative refinement of the queries themselves. The two are not exclusive: a future system could use both (queries refine via deformable cross-attention, memory refines via decoder feedback) and we expect the latter to compose with existing query-side improvements.

The gate-floor reparameterization is reminiscent of techniques used elsewhere when a learnable gate must be prevented from collapsing: highway networks~\cite{srivastava2015highway} use a similar additive bias term but for a different purpose (gradient flow), and various "soft skip-connection" formulations in transformer literature shape the residual branch's contribution explicitly. We are not aware of a prior application of a hard floor in the residual-attention setting; the design idea is small, but in our setting the difference between "unconstrained sigmoid" and "constrained sigmoid" is the difference between a silent and a measurable mechanism.
